% page 463 du fac-simile (image p-462 du scan)
% bloc 167.6 x 250.6 mm ; marge gauche 34.4 mm
\pgc{12.700mm}{0.000mm}{
\l{\cel{53}455}
\l{}
\l{\sou{postigo}. Vitro-plako di fenestro, quan onu povas apertar sen}
\l{\cel{3}apertar la tota fenestro, por parolar ad ulu, por lasar la ae-}
\l{\cel{3}ro enirar, e c. - S.}
\l{}
\l{\sou{postiliono}. I. La duktisto di posto-veturo, di qua la vesti esis}
\l{\cel{3}maxim-multa-kaze di koloro okul-frapanta, kun chapelo ruband-}
\l{\cel{3}oza, e qua kavalkis sur un de la kavali jungita. - II. Duesma}
\l{\cel{3}duktisto di karoso tirata da 4 o 6 kavali - la viro qua duktas}
\l{\cel{3}la kavali avana. - DEFIRS.}
\l{}
\l{\sou{postpoziciono}. (\sou{gram}.)Partikulo lokizita pos la vorto a qua lu}
\l{\cel{3}relatas. -}
\l{}
\l{\sou{postular}. (\sou{trans., de)}. I. Demandar rigoroze (ulo) konseque de}
\l{\cel{3}sua yuro, de sua autoritato, de sua forteso. - II. Igar ne-}
\l{\cel{3}kareebla. - FL.}
\l{}
\l{\sou{postulato}. Propoziciono quan onu prizentas kom egardenda,pro}
\l{\cel{3}lua exakteso, sen demonstrar lu. - DEFIRS.}
\l{}
\l{\sou{postuma}. Naskinta pos la morto di sua patro. (\sou{anke metaf.)}\cel{1}DEFIS.}
\l{}
\l{\sou{posturar}.\cel{2}(\sou{netrans}.) I. (\sou{aludante persono qua esas ava}n\cel{1}\sou{piktisto,}}
\l{\cel{3}\sou{fotografisto qui esas pronta operacar}). Movar sua korpo, sua}
\l{\cel{3}muskuli, modifikar la traiti di sua vizajo, quaze agas lo}
\l{\cel{3}modelo, segun la cirkonstanci. - II. Igar su aspektar plu}
\l{\cel{3}valoroza kam\cel{2}onu esas reale. - DEFIS.}
\l{}
\l{\sou{poto}. Menajo-vazo di qua la materio, la formo, la dimensioni}
\l{\cel{3}varias segun la uzi a qua lu esas destinita. - EFS.}
\l{}
\l{\sou{potaso}. (\sou{kemi}o).Karbonato di kalio quan onu extraktas de la}
\l{\cel{3}cindro de ula vejetanti. -\cel{1}\sou{Simbolo kemiala}\cel{1}: 2 KOH.( DEFIRS.}
\l{}
\l{\sou{potenco}. (\sou{mat}em.)\cel{2}Valoro di quanto nombrala, multiplikita un,}
\l{\cel{3}du o tri-foye per lu ipsa. - DEFIR.}
\l{}
\l{\sou{potencialo}. (\sou{mekan.}) Sumo de la forci qui esas necesa por startar}
\l{\cel{3}sistemo. - II. (\sou{fiziko)}\cel{2}Intenseso-grado di la forco necesa}
\l{\cel{3}por retroduktar molekulo deformita a lua formo prima.- DEFIS.}
\l{}
\l{\sou{potenciometro}. (\sou{elektr.}) Aparato, inventita da Clark, analoga a}
\l{\cel{3}la reostato de Wheatstone, uzata por mezurar la difero di pot-}
\l{\cel{3}enciali inter la poli di baterio.}
\l{}
\l{\sou{potenta}. I. Qua disponas la povo produktar granda efekti. -}
\l{\cel{3}II. Qua povas impozar sua autoritato. - DEFIS.}
\l{}
\l{\sou{povar.}\cel{2}(trans. e netrans.) Esar en la stando necesa por agar e}
\l{\cel{3}facar ulo, pro havar la forteso e la habileso necesa. - FIS.}
\l{}
\l{\sou{povra.}\cel{1}Qua ne havas lo necesa por suficar a su. (antonimo : richa)}
\l{\cel{3}- EFIS.}
\l{}
\l{pozar. (trans.)Situar en ta o ca plaso. - II. Situar en ca o}
\l{\cel{3}ta plaso, por permanar en ta loko. - FIS.}
}
